% !TEX root = ../main.tex
\clearpage
\appendix
\phantomsection
\addcontentsline{toc}{section}{Appendix}
\section*{Appendix}
\setcounter{figure}{0}
\renewcommand{\thefigure}{S\arabic{figure}}
\setcounter{table}{0}
\renewcommand{\thetable}{S\arabic{table}}

\phantomsection
\addcontentsline{toc}{subsection}{Appendix A. Metric Definitions}
\subsection*{Appendix A. Metric Definitions}

For each model evaluated on the screening benchmark, the following metrics are reported.
TP = true positives (included papers correctly retrieved);
FP = false positives (excluded papers incorrectly retrieved);
FN = false negatives (included papers missed);
TN = true negatives (excluded papers correctly excluded).
\[
\text{Recall} = \frac{\text{TP}}{\text{TP}+\text{FN}},
\quad
\text{Precision} = \frac{\text{TP}}{\text{TP}+\text{FP}},
\quad
F_1 = \frac{2 \times \text{Precision} \times \text{Recall}}{\text{Precision}+\text{Recall}}.
\]
\[
\text{NNS} = \frac{\text{TP}+\text{FP}}{\text{TP}},
\quad
\text{WR} = 1 - \frac{\text{TP}+\text{FP}}{\text{TP}+\text{FP}+\text{FN}+\text{TN}}.
\]
\[
\text{Excluded GT rate} = \frac{\text{FN}}{\text{FN}+\text{TN}}.
\]

\clearpage

\phantomsection
\addcontentsline{toc}{subsection}{Appendix B. Screening Model Comparison}
\subsection*{Appendix B. Screening Model Comparison}

Supplementary Figure~\ref{fig:model-compare} visualizes recall and papers needed to screen across candidate models and baselines. Table~\ref{tab:model-summary} reports the corresponding aggregate counts and metrics on the unified benchmark. The LEADS screening baseline follows the official LEADS title-and-abstract screening use case, using a disease- and parameter-specific prompt and returning a binary include/exclude decision. Prompt templates for the proposed workflow are provided in Appendix F.

\begin{figure}[H]
  \centering
  \includegraphics[width=\textwidth]{pdf/model_recall_nns_benchmark.pdf}
  \caption{Model-assisted screening performance across epidemiological parameters. a) Recall across three parameter domains, shown as dot plots to emphasize small differences among high-recall models. b) Papers needed to screen per included paper (\nns{}), where lower values indicate fewer papers screened per included paper; numbers above each panel indicate the original unassisted NNS.}
  \label{fig:model-compare}
\end{figure}

\papertable{model_summary.tex}

\clearpage

\phantomsection
\addcontentsline{toc}{subsection}{Appendix C. Second-Stage Full-Text Rescreening Analysis}
\subsection*{Appendix C. Second-Stage Full-Text Rescreening Analysis}

Tables~\ref{tab:fulltext-tier2-p13} and \ref{tab:p13-remaining-fp-audit} summarize a controlled comparison between title-and-abstract screening and full-text screening for the COVID-19 serial-interval review by Ali et al. 2022. The Ali et al. review was selected because it offered an appropriately sized raw pool and one of the largest sets of ground-truth included studies among the serial-interval tasks, providing sufficient signal to compare title-and-abstract and full-text inputs at the screening stage. Using the same five-criterion screening prompt and the same Strong, Possible, or Unlikely triage rule, screening was repeated with parsed full-text passages substituted for the title and abstract, while all other workflow components were held constant. Here, Stage II denotes this experimental full-text screening condition applied to records stratified by the Stage I title-and-abstract labels, not a required second step in the main screening pipeline.

\papertable{fulltext_tier2_p13_ablation.tex}

\papertable{fulltext_tier2_p13_remaining_fp_audit.tex}

\clearpage

\phantomsection
\addcontentsline{toc}{subsection}{Appendix D. Full-Text Parameter Coding and Extraction Comparison}
\subsection*{Appendix D. Full-Text Parameter Coding and Extraction Comparison}

Supplementary Figure~\ref{fig:supp-coding-leads} and Tables~\ref{tab:coding-leads-summary}--\ref{tab:coding-compare-leads-detail} provide the full per-source-review coding comparison between MetaAgent and LEADS. The detail table lists the source-review reference estimate used as the comparison target for each source review.

\begin{figure}[H]
  \centering
  \includegraphics[width=\textwidth]{pdf/coding_compare_leads_official_final.pdf}
  \caption{Supplementary coding comparison of source-review (SR), MetaAgent, and LEADS estimates across serial interval, reproduction number, and fatality.}
  \label{fig:supp-coding-leads}
\end{figure}

\papertable{coding_compare_leads_summary.tex}

\papertable{coding_compare_leads_detail.tex}

\clearpage

\phantomsection
\addcontentsline{toc}{subsection}{Appendix E. Benchmark Source-Review Inventory}
\subsection*{Appendix E. Benchmark Source-Review Inventory}

Table~\ref{tab:source-reviews} lists the source reviews used to construct the frozen screening and coding benchmark.

\papertable{source_reviews.tex}

\clearpage

\phantomsection
\addcontentsline{toc}{subsection}{Appendix F. Prompt Specifications}
\subsection*{Appendix F. Prompt Specifications}

The prompt templates below summarize the stable runtime structure used in the screening and full-text coding workflows. Project-specific disease names, parameter names, source-review questions, screening-profile guidance, codebook fields, and JSON schema constraints were filled programmatically at run time. Long parameter-specific examples and implementation-level formatting constraints are abbreviated for readability.

\begin{prompttemplate}{Template F1. Title-and-abstract screening}
\textbf{System role and task.} You are assisting with title-and-abstract screening for an epidemiological systematic review. Score the study on five relevance dimensions and classify it into one of three tiers using holistic judgment.

\textbf{Review context.} Use the following source-review fields:
\begin{promptitems}
  \item Research question: \promptvar{research\_question}.
  \item Disease focus and exclusions: \promptvar{disease\_focus}; \promptvar{disease\_exclude}.
  \item Target parameter and exclusions: \promptvar{parameter\_focus}; \promptvar{parameter\_exclude}.
\end{promptitems}

\textbf{Available evidence.} Use the title, abstract, publication year, record/create date, keywords, publication types, and MeSH terms. If the abstract is missing or sparse, avoid over-certainty and retain only records with concrete parameter or outcome cues.

\textbf{Assessment dimensions.} Score each dimension from 0 to 4 with a concise evidence-based rationale:
\begin{promptitems}
  \item Disease relevance: whether the paper investigates the target disease.
  \item Population relevance: whether the paper concerns humans or human-derived data.
  \item Location relevance: whether the study reports data from a real-world geographic setting relevant to the source review.
  \item Original empirical evidence: whether the paper reports original empirical data or primary analysis rather than only review, commentary, or theoretical content.
  \item Target-parameter relevance: whether the paper reports, estimates, or plausibly contains the target epidemiological parameter.
\end{promptitems}

\textbf{Parameter-specific guidance.} Apply the profile-specific guidance for the current disease and parameter domain. This guidance defined accepted synonyms, related quantities to distinguish, and common exclusion patterns; detailed edge-case examples are abbreviated here.

\textbf{Triage output.} Assign exactly one tier:
\begin{promptitems}
  \item \textbf{S}: strong candidate; clear evidence for the target disease, original empirical evidence, and target parameter.
  \item \textbf{P}: possible candidate; plausible but incomplete or ambiguous evidence requiring human review or full-text verification.
  \item \textbf{U}: unlikely candidate; no plausible match to one or more core eligibility criteria.
\end{promptitems}

\textbf{Response format.} Return JSON only. Include the five dimension scores and rationales, an optional overall score and confidence, the final S/P/U tier, and a short overall justification. The runtime prompt additionally supplied the JSON schema constraint.
\end{prompttemplate}

\begin{prompttemplate}{Template F2. Full-text evidence localization}
\textbf{System role and task.} You are an epidemiology research assistant creating a structured index from the full text of a primary study. Use only evidence from the provided text; do not guess or infer unsupported values. This is a high-recall indexing stage, not final extraction.

\textbf{Input.} Use \promptvar{title}, \promptvar{abstract}, parsed full text, recoverable tables and captions, and the parameter-domain codebook.

\textbf{Indexing targets.} Identify document sections, tables, figure captions, and paragraphs that may contain evidence for the target parameter, uncertainty interval, sample size, numerator or denominator, population, geography, study period, estimation method, or review-specific synthesis rule. Domain-specific prompts emphasized serial-interval, reproduction-number, or fatality evidence as appropriate.

\textbf{Response format.} Return JSON only. Summarize document structure, relevant tables, study metadata, parameter mentions, analysis summaries when applicable, and field-level evidence locations. Use null or empty arrays when evidence is not found.
\end{prompttemplate}

\begin{prompttemplate}{Template F3. Full-text parameter coding and extraction}
\textbf{System role and task.} You are an epidemiology meta-analyst extracting codebook-defined records from a full-text article. Use only evidence from the provided text and the Stage-A index; do not guess or compute values unless explicitly instructed by the codebook.

\textbf{Input.} Use \promptvar{title}, \promptvar{abstract}, parsed full text, recoverable table summaries, the Stage-A evidence index, and the parameter-specific codebook.

\textbf{Extraction rules.} Extract the study-level record or records needed for the source-review parameter analysis. Follow field definitions strictly; preserve reported point estimates, uncertainty intervals, units, denominators, populations, study periods, and evidence locations. Prefer the headline or overall estimate unless the source-review codebook specifies a subgroup or denominator. Parameter-domain rules distinguished serial interval from related timing quantities, $R_0$ from $R_t$ and assumed model inputs, and CFR from IFR, hospitalized fatality, ICU fatality, or denominator-specific fatality estimates.

\textbf{Response format.} Return a JSON array only. Each record includes the codebook-defined fields, including parameter type, point estimate, uncertainty interval when reported, population or subgroup, study period, geography, numerator or denominator where applicable, evidence location, and notes. Use null for missing numeric fields and ``NR'' for required text fields that are not reported. The runtime prompt also supplied codebook-specific field descriptions and an example JSON output.
\end{prompttemplate}
